\section*{Capítulo 4}

\solucion{4.1} (a) Corta: si el combustible sube, pagará más. (b) Larga en dólares: cobrará dólares, que valen más soles si el dólar sube. (c) Corta en soles: debe más soles de los que tiene. (d) Corto en dólares: si el dólar sube, pagará más soles.

\solucion{4.2} (a) Especulación: no tiene exposición previa. (b) Cobertura: reduce el riesgo de un inventario que ya posee. (c) Arbitraje: explota dos precios del mismo activo sin riesgo.

\solucion{4.3} Si el WTI sube US\$ 5, la posición corta pierde $5\times1{,}000=5{,}000$ por contrato. El saldo cae de 6,000 a 1,000, por debajo del mantenimiento, y debe reponerse hasta el margen inicial: 5,000 por contrato, $2{,}901\times5{,}000=14.5$ millones en efectivo. Si el WTI baja US\$ 5, la cuenta gana 5,000 por contrato y la empresa puede retirar ese excedente: no hay que depositar nada.

\solucion{4.4} $\text{CVA}\approx0.55\times25\times0.12=1.65$ millones.

\solucion{4.5} Porque le permite acumular pérdidas en el derivado (que se compensan con ganancias en la posición cubierta) sin entregar caja. El banco, en cambio, queda expuesto al riesgo de crédito de PETROPERÚ hasta ese umbral; con una calificación B-/Caa1 solo lo aceptará cobrando un CVA mayor o con la garantía de un tercero solvente, como el Estado.

\solucion{4.6} El payoff es el pago del contrato al vencimiento: 3. La ganancia descuenta la prima pagada: $3-5=-2$. La opción \enquote{pagó}, pero la operación perdió 2.

\solucion{4.7} Pasivo de 95, umbral 79: exceso de 16, mayor que el MTA de 5, así que entrega 16 millones. Con un pasivo de 82, el exceso es 3, menor que el MTA: no entrega nada (todavía).

\solucion{4.8} $\lambda\approx0.03/0.55=5.45\,\%$ anual. Probabilidad de incumplir en el primer año: $1-e^{-0.0545}=5.3\,\%$.

\solucion{4.9} Porque PETROPERÚ ya incumplió un covenant no financiero del CESCE (por la opinión con salvedades del auditor). Con un \emph{cross-default} estándar, ese incumplimiento podría permitir al banco terminar los derivados desde el primer día, justo cuando más se necesitan y posiblemente con un pago por terminación. Con \emph{cross-acceleration} solo cuenta si el acreedor del CESCE efectivamente acelera la deuda, y la dispensa expresa elimina ese riesgo para el CESCE.
